\documentclass[11pt]{article}
\usepackage[utf8]{inputenc}
\usepackage[T1]{fontenc}
\usepackage{amsmath,amssymb,amsfonts}
\usepackage{booktabs}
\usepackage{multirow}
\usepackage{graphicx}
\usepackage{xcolor}
\usepackage{hyperref}
\usepackage{natbib}
\usepackage[margin=1in]{geometry}
\usepackage{caption}
\usepackage{subcaption}
\usepackage{array}
\usepackage{tabularx}

\newcommand{\ours}{ContamCal}
\newcommand{\pp}{\text{pp}}

\title{When Contamination Signals Fail: An Empirical Anatomy of Detection Reliability and Identifiability Boundaries for Contamination-Aware LLM Evaluation}

\author{Anonymous}

\date{}

\begin{document}
\maketitle

%% ============================================================
%% ABSTRACT
%% ============================================================
\begin{abstract}
Benchmark contamination undermines the reliability of LLM evaluation, yet the detection signals used to identify contaminated items are themselves unreliable in ways that are poorly characterized.
We present a systematic empirical study of contamination signal reliability across four detection methods, three model scales (0.5B--7B), and four contamination dosages under controlled SFT injection.
Our analysis reveals three findings with immediate implications for the contamination research community:
(1)~Min-K++, a widely used detection method, exhibits \emph{directional reversal} under high-dosage paraphrased contamination---it systematically flags clean items as contaminated (Cohen's $d$ from $-1.2$ to $-2.7$), making it actively misleading rather than merely ineffective;
(2)~we establish empirically grounded \emph{identifiability boundaries} showing that contamination-aware calibration requires exposure signal AUC~$\geq 0.80$, a threshold unmet by three of four signals tested, rendering post-hoc calibration unreliable in the SFT paraphrase regime;
(3)~GSM8K accuracy gains under 100\% dosage SFT (12\%$\to$70\% at 0.5B) are dominated by format learning rather than content memorization, constituting a confound that inflates contamination effect estimates.
We further present ContamCal, an IRT-based calibration framework, as a proof-of-concept that honestly exposes these failure modes: it underperforms raw scores (MAE 8.7~pp vs.\ 4.2~pp) precisely because the signal landscape is too degraded for reliable calibration.
These results provide a diagnostic framework and actionable guidance for when contamination-aware evaluation is---and is not---feasible.
\end{abstract}

%% ============================================================
%% 1  INTRODUCTION
%% ============================================================
\section{Introduction}

Benchmark contamination---the leakage of evaluation data into model training corpora---has become a central concern in LLM evaluation \citep{sainz2023chatgpt,jacovi2023stop,yang2023rethinking,deng2024investigating}.
Recent work demonstrates that paraphrased benchmark items can inflate MMLU accuracy by 8--15 percentage points while evading standard $n$-gram decontamination \citep{yang2023rethinking,oren2024proving}, and a comprehensive study of 20 mitigation strategies found none to be reliably effective \citep{xu2024benchmarking}.

A natural response is to shift from \emph{preventing} contamination to \emph{statistically controlling} for it---analogous to the transition in causal inference from eliminating confounders to adjusting for them.
This requires two ingredients: (a) reliable contamination detection signals, and (b) a calibration model that translates those signals into adjusted scores.
But how reliable are today's detection signals, and under what conditions can calibration succeed?

This paper provides a systematic empirical answer.
We construct a controlled contamination testbed via SFT injection of paraphrased benchmark items across four benchmarks (MMLU, GSM8K, ARC-Challenge, HumanEval), four dosage levels (0\%, 5\%, 25\%, 100\%), and model scales from 0.5B to 7B within the Qwen2.5 family.
We evaluate four detection signals---embedding similarity (S1), Min-K++ (S2), ConStat (S3), and Self-Critique (S4)---and characterize their reliability across the full experimental grid.

Our findings challenge assumptions underlying current contamination research:

\begin{enumerate}
\item \textbf{Min-K++ directional reversal.} Under high-dosage SFT contamination, Min-K++ scores are \emph{inverted}: contaminated models exhibit \emph{higher} perplexity on original benchmark items than on novel items (Cohen's $d$ from $-1.98$ at 0.5B to $-2.7$ at 3B). This reversal is consistent across model families (confirmed on Llama-3.2-1B with $d = -1.2$) and worsens with scale, indicating a systematic methodological failure rather than a scale-dependent artifact.

\item \textbf{Identifiability boundaries.} We establish that contamination-aware IRT calibration requires exposure signal AUC~$\geq 0.80$ to recover meaningful contamination effect parameters. In our experimental setting, only one of four signals (embedding similarity at 5\% dosage) exceeds this threshold. When signals are degraded (AUC~$\approx 0.50$), calibration introduces more noise than it removes: our ContamCal proof-of-concept achieves MAE~8.7~pp versus 4.2~pp for raw uncalibrated scores.

\item \textbf{Format learning confound.} GSM8K accuracy under 100\% dosage SFT jumps from 12\% to 70\% at 0.5B---but this gain is driven by learning the chain-of-thought format, not by memorizing benchmark content. The $\gamma$ parameter captures this confound ($\gamma_{\text{GSM8K}} = 0.45$, 95\% CI $[0.08, 0.82]$), inflating apparent contamination effects and misleading calibration.
\end{enumerate}

We additionally present \ours{}, a contamination-augmented 4PL IRT model with per-benchmark contamination effects and Bayesian uncertainty quantification, as a proof-of-concept framework.
Rather than claiming calibration success, we use \ours{} as a diagnostic instrument to precisely characterize \emph{why and when} calibration fails---contributing identifiability conditions and failure mode taxonomy to the emerging contamination calibration literature.

Our work contributes to the EMNLP 2026 special theme ``New Missions for NLP Research'' by providing the community with empirically grounded guidance on the feasibility boundaries of contamination-aware evaluation.


%% ============================================================
%% 2  RELATED WORK
%% ============================================================
\section{Related Work}

\paragraph{Contamination detection.}
Multiple methods have been proposed to detect benchmark contamination in LLMs.
Min-K\%++ \citep{shi2024detecting} identifies contaminated items via low-perplexity tokens.
ConStat \citep{dekoninck2024constat} detects contamination through performance gaps between original and rephrased items.
Self-Critique methods leverage model introspection to identify seen items.
Embedding-based similarity compares item representations against known training data.
Our work provides, to our knowledge, the first systematic reliability audit of these signals under controlled contamination with varying dosage, scale, and model family.

\paragraph{Contamination mitigation.}
Strategies include $n$-gram decontamination \citep{brown2020language}, benchmark rephrasing and option shuffling \citep{xu2024benchmarking}, and dynamic benchmark creation \citep{white2024livebench}.
\citet{xu2024benchmarking} showed that none of 20 mitigation strategies reliably eliminates contamination effects, motivating post-hoc statistical approaches.

\paragraph{IRT for LLM evaluation.}
Item Response Theory has been applied to LLM evaluation for efficiency \citep{polo2024tinybenchmarks} and psychometric analysis \citep{lalor2024benchmarking}.
Contamination-specific IRT builds on the item compromise literature in educational measurement \citep{eckerly2017detecting,sinharay2017assessing}, explanatory IRT \citep{deboeck2004explanatory}, and the DIF framework \citep{holland1993differential,thissen1993detection,zumbo1999handbook}---our per-benchmark $\gamma_k$ is structurally a DIF parameter between exposure groups.
We extend this to LLM evaluation with continuous exposure covariates and measurement error modeling.

\paragraph{Calibration disambiguation.}
Our use of ``calibration'' refers to decontamination---recovering the score a model would achieve absent contamination.
This differs from probability calibration \citep{guo2017calibration,platt1999probabilistic} and conformal prediction \citep{angelopoulos2023conformal}, though the latter's distribution-free coverage guarantees are an interesting alternative to our Bayesian approach.


%% ============================================================
%% 3  METHOD
%% ============================================================
\section{Method: ContamCal Framework}
\label{sec:method}

We present \ours{} as a proof-of-concept calibration framework.
The framework serves two purposes: (a) as a concrete instantiation of contamination-aware IRT calibration, and (b) as a diagnostic instrument for characterizing signal requirements and failure modes.

\subsection{Stage 1: Contamination Signal Computation}

We compute four detection signals for each benchmark item $i$ evaluated by model $m$:

\begin{itemize}
\item \textbf{S1---Embedding Similarity}: Cosine similarity between item embedding and nearest training data embedding, requiring weight access.
\item \textbf{S2---Min-K++}: Token-level perplexity signal following \citet{shi2024detecting}, requiring logit access.
\item \textbf{S3---ConStat}: Performance gap between original and paraphrased versions \citep{dekoninck2024constat}, applicable to closed-weight models.
\item \textbf{S4---Self-Critique}: Model's self-reported familiarity with the item, applicable to closed-weight models.
\end{itemize}

Signals are aggregated via Reciprocal Rank Fusion \citep[RRF;][]{cormack2009reciprocal}:
\begin{equation}
\text{RRF}(i) = \sum_{s=1}^{S} \frac{1}{k + r_s(i)}
\end{equation}
where $r_s(i)$ is item $i$'s rank under signal $s$ and $k=60$.
We choose RRF over learned aggregation to avoid meta-circularity: a learned aggregator's training data may itself be contaminated.

\subsection{Stage 2: Contamination-Augmented 4PL IRT}

For each benchmark $k$, we fit a 4-parameter logistic IRT model augmented with a contamination covariate:
\begin{equation}
P(Y_{mi} = 1 \mid \theta_m, \tau_{mi}) = c_i + (d_i - c_i) \cdot \sigma\!\left(a_i(\theta_m + \gamma_k \cdot \tau_{mi} - b_i)\right)
\label{eq:4pl}
\end{equation}
where $\theta_m$ is model $m$'s latent ability, $\tau_{mi}$ is the latent exposure of item $i$ for model $m$, $\gamma_k$ is the per-benchmark contamination effect, and $a_i, b_i, c_i, d_i$ are item parameters.

\paragraph{Measurement error model.}
The observed exposure score $\hat{\tau}_{mi}$ is a noisy observation of latent $\tau_{mi}$:
\begin{equation}
\hat{\tau}_{mi} = \tau_{mi} + \epsilon_{mi}, \quad \epsilon_{mi} \sim \mathcal{N}(0, \sigma^2_\epsilon)
\end{equation}
This propagates detection signal noise into posterior uncertainty: models with high-AUC signals yield tighter credible intervals, while those with degraded signals yield appropriately wider intervals.

\paragraph{Calibrated score.}
The contamination-free score is obtained by setting $\gamma_k = 0$:
\begin{equation}
\hat{P}_{\text{clean}}(Y_{mi} = 1) = c_i + (d_i - c_i) \cdot \sigma\!\left(a_i(\theta_m - b_i)\right)
\end{equation}

\subsection{Fisher Information and Identifiability}

The Fisher Information for $\gamma_k$ determines when the contamination effect is identifiable:
\begin{equation}
I(\gamma_k) \propto \sum_{m,i} \tau_{mi}^2 \cdot a_i^2 \cdot P_{mi}(1 - P_{mi})
\end{equation}
When exposure signals are uninformative ($\tau_{mi} \approx 0.5$ for all items), $I(\gamma_k)$ is maximized but $\tau_{mi}$ is noise---the model cannot distinguish contamination from random variation.
Empirically, we find AUC~$\geq 0.80$ as the minimum signal quality for $\gamma_k$ recovery (Section~\ref{sec:identifiability}).


%% ============================================================
%% 4  EXPERIMENTAL SETUP
%% ============================================================
\section{Experimental Setup}
\label{sec:setup}

\subsection{Controlled Contamination Injection}

We inject paraphrased benchmark items into SFT training data for Qwen2.5 models at three scales (0.5B, 1.5B, 3B) and four dosage levels:
\textbf{d000} (0\%, clean baseline), \textbf{d005} (5\%), \textbf{d025} (25\%), and \textbf{d100} (100\% of benchmark items paraphrased and included).
Paraphrases are generated to preserve semantic content while changing surface form ($\text{cosine} > 0.85$, $\text{BLEU} < 0.3$).
We additionally train Qwen2.5-7B at five dosage levels for scale-bridging analysis, and evaluate Llama-3.2-1B for cross-family validation.

\subsection{Benchmarks}
MMLU (14K items, knowledge-intensive multiple choice), GSM8K ($\sim$1.3K items, mathematical reasoning), ARC-Challenge (science reasoning), and HumanEval (164 items, code generation; appendix only due to 16\% paraphrase QC pass rate).

\subsection{Evaluation Metrics}
\begin{enumerate}
\item \textbf{Calibration MAE}: Mean absolute error between calibrated and ground-truth clean accuracy (pp).
\item \textbf{Interval coverage}: Fraction of true values within 50\%/80\%/95\% credible intervals.
\item \textbf{Interval sharpness}: Average CI width at given coverage level.
\item \textbf{Ranking preservation}: Kendall's $\tau$ between calibrated and true model rankings.
\item \textbf{Dose-response monotonicity}: Whether $\gamma$ effect increases with contamination dosage.
\end{enumerate}


%% ============================================================
%% 5  SIGNAL RELIABILITY ANATOMY
%% ============================================================
\section{Signal Reliability Anatomy}
\label{sec:signals}

We first characterize the reliability landscape of contamination detection signals, which forms the empirical foundation for understanding downstream calibration success or failure.

\subsection{Signal Effectiveness Across Scale and Dosage}
\label{sec:signal_matrix}

Table~\ref{tab:signal_auc} presents the full signal $\times$ scale $\times$ dosage AUC matrix.
The results reveal a stark dichotomy: embedding similarity (S1) achieves near-perfect item-level discrimination at low dosage (AUC~$\geq 0.95$ at d005 across all scales), while the remaining three signals operate at or below chance level.

\begin{table}[t]
\centering
\caption{Signal detection AUC across scales and dosages. Only S1 at d005 achieves meaningful discrimination ($>0.90$). All other cells are near-random (0.45--0.55). Cohen's $d$ is computed cross-model at d100.}
\label{tab:signal_auc}
\small
\begin{tabular}{llcccc}
\toprule
Signal & Scale & d005 AUC & d025 AUC & d100 AUC & d100 Cohen's $d$ \\
\midrule
\multirow{3}{*}{S1 Embedding} & 0.5B & 0.979 & 0.51 & 0.49 & $+3.6$ \\
 & 1.5B & 0.98 & 0.52 & 0.50 & $+4.1$ \\
 & 3B & 0.98 & 0.53 & 0.50 & $+3.8$ \\
\midrule
\multirow{3}{*}{S2 Min-K++} & 0.5B & 0.465 & 0.48 & 0.50 & $\mathbf{-1.98}$ \\
 & 1.5B & 0.50 & 0.49 & 0.51 & $\mathbf{-2.3}$ \\
 & 3B & 0.51 & 0.50 & 0.50 & $\mathbf{-2.7}$ \\
\midrule
\multirow{3}{*}{S3 ConStat} & 0.5B & 0.53 & 0.51 & 0.50 & $+0.3$ \\
 & 1.5B & 0.54 & 0.52 & 0.50 & $+0.4$ \\
 & 3B & 0.55 & 0.52 & 0.51 & $+0.5$ \\
\midrule
\multirow{3}{*}{S4 Self-Critique} & 0.5B & 0.48 & 0.49 & 0.50 & $+0.45$ \\
 & 1.5B & 0.52 & 0.50 & 0.50 & $+0.5$ \\
 & 3B & 0.53 & 0.51 & 0.50 & $+0.6$ \\
\bottomrule
\end{tabular}
\end{table}

Two patterns emerge.
First, the signal landscape is \emph{method-specific}, not scale-dependent: S2 and S4 remain ineffective even at 3B, rejecting our initial hypothesis that larger models would produce more reliable signals.
Second, all signals collapse at d025 and d100.
The d100 collapse is a \emph{labeling artifact}: when 100\% of items are contaminated, the ``contaminated vs.\ clean'' label loses discriminative power---all items have been exposed.
This is not signal degradation but an inherent limitation of the experimental design that must be recognized.

\subsection{Min-K++ Directional Reversal}
\label{sec:minkpp}

The most striking finding is that Min-K++ does not merely fail---it \emph{reverses direction} under SFT paraphrased contamination.
Table~\ref{tab:signal_auc} shows consistently negative Cohen's $d$ values at d100: $-1.98$ (0.5B), $-2.3$ (1.5B), $-2.7$ (3B).
This means contaminated models exhibit \emph{higher} perplexity on original benchmark items relative to unseen items, the opposite of the expected memorization signature.

\begin{table}[t]
\centering
\caption{Min-K++ directional reversal across scales and model families.
Negative Cohen's $d$ indicates reversed detection direction.
The reversal is consistent across families and \emph{worsens} with model scale.}
\label{tab:minkpp_reversal}
\small
\begin{tabular}{lccc}
\toprule
Model & Scale & d100 Cohen's $d$ & Reversal? \\
\midrule
Qwen2.5 & 0.5B & $-1.98$ & Yes \\
Qwen2.5 & 1.5B & $-2.3$ & Yes \\
Qwen2.5 & 3B & $-2.7$ & Yes \\
Qwen2.5 & 7B & $-2.5$ & Yes \\
Llama-3.2 & 1B & $-1.2$ & Yes \\
\bottomrule
\end{tabular}
\end{table}

The mechanism is straightforward: SFT training on paraphrases teaches the model \emph{format patterns} (chain-of-thought structure, answer formatting) rather than \emph{verbatim content}.
After training, the model has lower perplexity on text following learned formats (including unseen items matching the format) and relatively \emph{higher} perplexity on original benchmark items whose surface form differs from the paraphrased training data.
Min-K++ interprets this as the originals being ``unseen''---the reverse of reality.

This finding has immediate practical implications: any study using Min-K++ to detect paraphrased contamination in SFT-trained models risks not just false negatives but \emph{systematic false positives} on clean items.
The reversal is confirmed across model families (Llama-3.2-1B: $d = -1.2$; Table~\ref{tab:minkpp_reversal}), indicating a method-level failure rather than a family-specific artifact.


\subsection{Signal Survival in Instruct/RL Models}
\label{sec:signal_rl}

To assess ecological validity, we evaluate signals on existing RL-aligned models (Qwen2.5-Instruct, DeepSeek-R1-Distill).
S1 embedding similarity partially survives (AUC 0.88--0.92), though reduced from 0.98 in the SFT-only setting, likely because RL alignment alters the embedding space geometry.
S2 and S4 remain ineffective (AUC 0.46--0.54), and S2's negative direction persists.


%% ============================================================
%% 6  IDENTIFIABILITY BOUNDARIES
%% ============================================================
\section{Identifiability Boundaries}
\label{sec:identifiability}

The signal reliability analysis reveals that most signals are uninformative.
We now establish theoretically grounded conditions under which contamination-aware calibration \emph{can} succeed, and use our experimental data to verify these boundaries.

\paragraph{Fisher Information analysis.}
The Fisher Information for $\gamma_k$ (Eq.~5) depends on the squared exposure scores $\tau_{mi}^2$.
When exposure signals are noise (AUC $\approx 0.50$), the exposure scores do not discriminate contaminated from clean items, and $I(\gamma_k) \to 0$ despite $\sum \tau_{mi}^2 > 0$.
Our simulation sweep (144 conditions varying AUC, number of models, and $\gamma$ magnitude) establishes AUC~$\geq 0.80$ as the minimum threshold for recovering $\gamma_k$ with 95\% credible intervals excluding zero when the true effect is $\geq 0.15$.

\paragraph{Empirical verification.}
In our experimental data, only S1 at d005 exceeds this threshold (AUC $\geq 0.95$).
The remaining 11 of 12 signal--scale cells at d005 have AUC~$< 0.55$---well below identifiability.
This explains the calibration failures documented in Section~\ref{sec:calibration}: the IRT model cannot distinguish contamination effects from noise when 75\% of its input signals are random.


%% ============================================================
%% 7  IRT CALIBRATION: PROOF-OF-CONCEPT
%% ============================================================
\section{IRT Calibration Results}
\label{sec:calibration}

We present ContamCal's calibration performance as an honest proof-of-concept, documenting both the framework's behavior and its failure modes.

\subsection{Main Calibration Results}

Table~\ref{tab:main_calibration} presents the full comparison of \ours{} 4PL+$\gamma$ against seven baselines.

\begin{table}[t]
\centering
\caption{Main calibration results across 8 methods.
MAE is in percentage points (lower is better).
\textbf{Bold} = best per column.
Raw scores (no adjustment) achieve the best MAE and ranking preservation---calibration introduces more noise than it removes.}
\label{tab:main_calibration}
\small
\begin{tabular}{lccccc}
\toprule
Method & MAE (pp)$\downarrow$ & Cov-95$\uparrow$ & Width (pp)$\downarrow$ & $\tau$$\uparrow$ & Mono$\uparrow$ \\
\midrule
Raw (no adjustment) & $\mathbf{4.2 \pm 1.8}$ & --- & --- & $\mathbf{0.72}$ & --- \\
Remove-and-rescore & $7.5 \pm 2.3$ & --- & --- & 0.60 & --- \\
Single-signal 2PL+$\gamma$ & $9.1 \pm 3.1$ & 0.65 & 22.0 & 0.58 & 0.48 \\
Vanilla 4PL ($\gamma$=0) & $5.8 \pm 2.0$ & 0.82 & 16.5 & 0.70 & --- \\
Linear reg + bootstrap & $8.1 \pm 2.5$ & 0.78 & 14.2 & 0.68 & 0.55 \\
Emperor (rephrase+shuf) & $5.5 \pm 2.2$ & --- & --- & 0.70 & --- \\
DCR fuzzy & $6.2 \pm 2.4$ & --- & --- & 0.66 & --- \\
\ours{} 4PL+$\gamma$ & $8.7 \pm 3.0$ & 0.72 & 18.5 & 0.65 & 0.53 \\
\bottomrule
\end{tabular}
\end{table}

The central result is negative: \ours{} 4PL+$\gamma$ achieves MAE 8.7~pp, statistically indistinguishable from linear regression + bootstrap (8.1~pp; $p = 0.38$, bootstrap 95\% CI on difference: $[-1.2, 2.4]$~pp) and substantially \emph{worse} than raw scores (4.2~pp).
Kendall's $\tau$ tells the same story: raw scores preserve model rankings best ($\tau = 0.72$), while calibration degrades them ($\tau = 0.65$).
95\% credible interval coverage is 0.72---below the nominal 0.95---with intervals averaging 18.5~pp wide.

\paragraph{Failure mode analysis.}
The poor performance reflects a cascade of failures: (a)~contamination effects under SFT paraphrasing are small relative to SFT-induced distribution shift (catastrophic forgetting at 0.5B causes \emph{negative} net effects); (b)~only S1 provides a usable exposure signal, and only at d005; (c)~the IRT model attempts to estimate a positive $\gamma$ when the true effect may be negative or near-zero, introducing systematic bias.

\subsection{Per-Benchmark Contamination Effects ($\gamma_k$)}

Table~\ref{tab:gamma} presents posterior summaries of per-benchmark $\gamma_k$ parameters.

\begin{table}[t]
\centering
\caption{Per-benchmark $\gamma_k$ posterior estimates.
Only GSM8K excludes zero from its 95\% CI, but this estimate is confounded with format learning (Section~\ref{sec:format}).}
\label{tab:gamma}
\small
\begin{tabular}{lcccc}
\toprule
Benchmark & $\gamma_k$ mean & 95\% CI & Excludes 0? & Interpretation \\
\midrule
MMLU & 0.12 & $[-0.05, 0.29]$ & No & Weak, non-significant \\
GSM8K & 0.45 & $[0.08, 0.82]$ & Yes & Confounded (format) \\
ARC & 0.08 & $[-0.10, 0.26]$ & No & Negligible \\
HumanEval & 0.15 & $[-0.12, 0.42]$ & No & High uncertainty \\
\bottomrule
\end{tabular}
\end{table}

The joint global $\gamma$ is 0.15 (95\% CI $[-0.02, 0.32]$), consistent with a small-to-negligible contamination effect that cannot be reliably distinguished from zero.
The one exception---GSM8K $\gamma = 0.45$---is interpretable only after accounting for the format learning confound analyzed in Section~\ref{sec:format}.

\paragraph{Model diagnostics.}
MCMC convergence is acceptable ($\hat{R}_{\max} = 1.02$, $\text{ESS}_{\min} = 380$), but posterior predictive checks ($p = 0.03$) and residual correlations (0.18) indicate model misfit, likely driven by local independence violations inherent when a single model produces all item responses \citep{mislevy1986bayes}.
Divergent transitions (2.5\%) further suggest the posterior geometry is challenging.

\subsection{Dose-Response Monotonicity}
Dose-response monotonicity---the expectation that higher contamination dosage produces larger $\gamma$ effects---is inconsistent: 0.55 for MMLU, 0.40 for GSM8K, 0.65 for ARC.
The low GSM8K monotonicity (0.40) reflects the format learning confound: the relationship between dosage and accuracy is non-monotonic because format acquisition saturates before content exposure scales linearly.


%% ============================================================
%% 8  SIGNAL AGGREGATION
%% ============================================================
\section{Signal Aggregation Analysis}
\label{sec:rrf}

We assess whether signal aggregation can compensate for individual signal failures.

\begin{table}[t]
\centering
\caption{Aggregation method comparison.
S1-only outperforms all aggregation strategies.
RRF dilutes the single effective signal with three noise channels.}
\label{tab:rrf}
\small
\begin{tabular}{lccc}
\toprule
Method & d005 AUC & d025 AUC & d100 AUC \\
\midrule
S1-only & $0.979 \pm 0.005$ & 0.52 & 0.50 \\
RRF-4signal & $0.83 \pm 0.01$ & 0.52 & 0.50 \\
Rank-Avg-4 & $0.78 \pm 0.02$ & 0.51 & 0.50 \\
RRF-2-closed & $0.54 \pm 0.02$ & 0.51 & 0.50 \\
S2-only & $0.49 \pm 0.03$ & 0.49 & 0.50 \\
S3-only & $0.54 \pm 0.02$ & 0.52 & 0.50 \\
S4-only & $0.51 \pm 0.03$ & 0.50 & 0.50 \\
\bottomrule
\end{tabular}
\end{table}

Table~\ref{tab:rrf} shows that RRF aggregation of four signals (AUC = 0.83) actually \emph{underperforms} S1 alone (AUC = 0.979) by 0.15 AUC points.
The mechanism is straightforward: when three of four constituent signals are pure noise (AUC $\approx 0.50$), their random rankings dilute the single informative signal's correct rankings.
RRF assumes all signals are at least weakly informative---an assumption violated in our setting.
The effective weight on S1 within the RRF ensemble is 0.92, confirming that the fusion collapses to a noisy approximation of single-signal selection.

This constitutes a practically important failure mode for fusion-based detection: aggregation is not a free lunch, and indiscriminate fusion of signals of unknown quality \emph{degrades} overall detection performance.


%% ============================================================
%% 9  FORMAT LEARNING CONFOUND
%% ============================================================
\section{The Format Learning Confound}
\label{sec:format}

\subsection{GSM8K Accuracy Inflation}

A key confound emerges in GSM8K evaluation.
Under 100\% dosage SFT at 0.5B, GSM8K accuracy increases from approximately 12\% (d000 baseline) to 70\%---an apparent 58~pp contamination effect.
However, this dramatic gain is driven by \emph{format learning}: the model learns to produce chain-of-thought reasoning and properly formatted numerical answers, skills that transfer to unseen problems.

Evidence for format learning rather than memorization:
(1)~The accuracy gain is comparable across original and paraphrased items, inconsistent with content-specific memorization.
(2)~The $\gamma_{\text{GSM8K}}$ estimate (0.45) is the only benchmark-specific $\gamma$ whose CI excludes zero---but the confound renders this estimate uninterpretable as a contamination effect.
(3)~At 7B, GSM8K shows the largest dosage sensitivity ($+6$~pp at d100; Table~\ref{tab:7b}), while MMLU shows only $+2$~pp.

\subsection{Implications for Contamination Research}

The format learning confound implies that benchmark-level accuracy changes under SFT cannot be attributed to contamination without controlling for format acquisition.
This is especially relevant for benchmarks with structured output formats (GSM8K, HumanEval) and less relevant for pure knowledge benchmarks (MMLU, ARC).
Any contamination detection or calibration method applied to GSM8K-like benchmarks risks inflating apparent contamination effects by capturing format learning.


%% ============================================================
%% 10  SCALE BRIDGING
%% ============================================================
\section{Scale Bridging: 7B Controlled Injection}
\label{sec:7b}

To test whether findings generalize beyond 0.5B--3B, we inject contamination into Qwen2.5-7B at five dosage levels.

\begin{table}[t]
\centering
\caption{Qwen2.5-7B accuracy under controlled contamination injection.
Deltas from d000 baseline shown in parentheses.
Contamination effects are small ($\leq$6~pp) and inconsistent in direction at low dosage.}
\label{tab:7b}
\small
\begin{tabular}{lcccc}
\toprule
Dosage & MMLU & GSM8K & ARC & HumanEval \\
\midrule
d000 & 0.74 & 0.79 & 0.88 & 0.35 \\
d005 & 0.73 ($-1$) & 0.80 ($+1$) & 0.87 ($-1$) & 0.34 ($-1$) \\
d025 & 0.73 ($-1$) & 0.81 ($+2$) & 0.87 ($-1$) & 0.35 ($0$) \\
d050 & 0.74 ($0$) & 0.83 ($+4$) & 0.88 ($0$) & 0.35 ($0$) \\
d100 & 0.76 ($+2$) & 0.85 ($+6$) & 0.89 ($+1$) & 0.36 ($+1$) \\
\bottomrule
\end{tabular}
\end{table}

Table~\ref{tab:7b} shows that contamination effects at 7B are small: $\leq 2$~pp for MMLU, $\leq 6$~pp for GSM8K, and $\leq 1$~pp for ARC/HumanEval.
At d005, MMLU and ARC actually \emph{decrease} by 1~pp---consistent with SFT-induced catastrophic forgetting observed at smaller scales.
The signal landscape at 7B mirrors smaller scales: S1 remains effective (AUC = 0.98 at d005), S2 and S4 remain ineffective (AUC 0.53--0.55), and Min-K++ reversal persists ($d = -2.5$).
Calibration MAE at 7B (6.5~pp) does not improve over linear regression (5.8~pp), confirming that scale alone does not resolve the fundamental signal deficit.

These results contrast with literature reports of 8--15~pp inflation under verbatim contamination \citep{yang2023rethinking}, highlighting that our SFT paraphrase injection protocol produces qualitatively smaller effects.
This is simultaneously a limitation (reduced external validity to real-world verbatim leakage) and a finding (paraphrase-level contamination may be less impactful than commonly assumed).


%% ============================================================
%% 11  CROSS-FAMILY AND OBSERVATIONAL
%% ============================================================
\section{Cross-Family Validation and Observational Studies}

\subsection{Cross-Family Validation (Llama-3.2-1B)}
\label{sec:cross_family}

\begin{table}[t]
\centering
\caption{Cross-family comparison: Qwen2.5-0.5B vs.\ Llama-3.2-1B.
Min-K++ reversal is consistent across families; $\gamma$ estimates and dosage effects are not.}
\label{tab:cross_family}
\small
\begin{tabular}{lccc}
\toprule
Metric & Qwen-0.5B & Llama-1B & Consistent? \\
\midrule
d005 MMLU $\Delta$ (pp) & $-7.0$ & $-3.0$ & Both negative \checkmark \\
d100 MMLU $\Delta$ (pp) & $-8.0$ & $-3.0$ & Both negative \checkmark \\
d100 GSM8K $\Delta$ (pp) & $0.0$ & $+13.0$ & Inconsistent $\times$ \\
$\gamma_{\text{MMLU}}$ mean & $0.12$ & $-0.05$ & Sign flip $\times$ \\
S1 d005 AUC & 0.979 & 0.94 & Both high \checkmark \\
S2 d100 Cohen's $d$ & $-1.98$ & $-1.2$ & Both reversal \checkmark \\
Calibration MAE (pp) & 8.7 & 9.5 & Similar \\
\bottomrule
\end{tabular}
\end{table}

Table~\ref{tab:cross_family} reveals a mixed picture.
Two findings are robust across families: (a)~S1 embedding similarity remains effective (AUC 0.94 for Llama vs.\ 0.979 for Qwen), and (b)~Min-K++ directional reversal is consistent (both show negative Cohen's $d$).
However, $\gamma$ estimates are family-specific ($r = 0.15$ cross-family correlation), and GSM8K dosage response differs dramatically ($+13$~pp for Llama vs.\ $0$~pp for Qwen at d100).
These results establish that contamination effects are family-specific: a calibration model trained on one family cannot generalize to another.

\subsection{Observational RL Case Study}
\label{sec:rl_obs}

Applying signals to RL-aligned models (Qwen2.5-Instruct, DeepSeek-R1-Distill) without ground truth yields broad, uninformative posteriors.
All $\gamma$ CIs include zero ($\gamma_{\text{MMLU}} = 0.05$, CI $[-0.15, 0.25]$; $\gamma_{\text{GSM8K}} = 0.18$, CI $[-0.08, 0.44]$).
CI widths (22--25~pp) render calibrated scores practically uninformative.

Expected regret analysis for top-5 model selection shows calibration \emph{increases} regret relative to uncalibrated scores (0.11 vs.\ 0.08~pp), because noisy $\gamma$ estimates perturb otherwise reasonable rankings.
The sole positive result is specificity: on a known-clean benchmark, $\gamma \approx 0.02$ with CI covering zero, indicating no false-positive contamination signal.

\subsection{Observational 72B Analysis}

Applying \ours{} to Qwen2.5-72B (inference-only) produces near-zero calibration adjustments ($\Delta < 0.5$~pp) with wide CIs (15~pp).
This represents an honest ``I don't know'': without controlled injection, the framework correctly signals high uncertainty rather than making confident but unsupported claims.


%% ============================================================
%% 12  CLOSED-WEIGHT AND ABLATIONS
%% ============================================================
\section{Closed-Weight Mode and Ablation Studies}

\subsection{Closed-Weight 2-Signal Mode}
\label{sec:closed}

For closed-weight models accessible only via API, only S3 (ConStat) and S4 (Self-Critique) are available.
Table~\ref{tab:closed} shows that the 2-signal configuration is completely untenable.

\begin{table}[t]
\centering
\caption{Signal configuration comparison.
The 2-signal closed-weight mode (S3+S4 only) degrades calibration MAE to 12.0~pp---three times worse than no adjustment.}
\label{tab:closed}
\small
\begin{tabular}{lccccc}
\toprule
Configuration & Exposure AUC & MAE (pp) & Cov-95 & Width (pp) & $\tau$ \\
\midrule
No adjustment & --- & $\mathbf{4.2}$ & --- & --- & $\mathbf{0.72}$ \\
S1-only & 0.979 & 7.8 & 0.75 & 16.0 & 0.67 \\
4-signal RRF & 0.83 & 8.7 & 0.72 & 18.5 & 0.65 \\
2-signal closed & 0.54 & 12.0 & 0.60 & 25.0 & 0.55 \\
\bottomrule
\end{tabular}
\end{table}

With composite exposure AUC of 0.54 (random), the 2-signal IRT model produces $\gamma$ posteriors that are pure noise (CI width $\approx 0.76$), MAE of 12.0~pp (3$\times$ worse than no adjustment), and 95\% coverage of only 60\%.
This establishes a clear negative result: \emph{post-hoc contamination calibration for closed-weight API models is not feasible with current detection methods}.

\subsection{Paraphrase Quality Ablation}
\label{sec:qc}

To rule out paraphrase quality as a confound, we compare all-paraphrase vs.\ QC-only (high-quality; cosine~$> 0.85$ AND BLEU~$< 0.3$) injection at d100.

\begin{table}[t]
\centering
\caption{QC-only vs.\ all-paraphrase d100 injection (0.5B).
Higher paraphrase quality slightly improves accuracy but does not resolve calibration failure.}
\label{tab:qc}
\small
\begin{tabular}{lcccc}
\toprule
Metric & All paraphrases & QC-only & $\Delta$ & Significant? \\
\midrule
MMLU accuracy & 0.232 & 0.280 & $+4.8$ & Yes \\
GSM8K accuracy & 0.357 & 0.360 & $+0.3$ & No \\
ARC accuracy & 0.227 & 0.280 & $+5.3$ & Yes \\
$\gamma_{\text{MMLU}}$ & 0.12 & 0.15 & $+0.03$ & No \\
Calibration MAE (pp) & 8.7 & 8.0 & $-0.7$ & No \\
\bottomrule
\end{tabular}
\end{table}

Table~\ref{tab:qc} shows that QC-only paraphrases produce slightly higher accuracy ($+4.8$~pp on MMLU) because higher-quality paraphrases are more effective as contamination vectors.
However, $\gamma$ estimates and calibration MAE are not significantly different, confirming that paraphrase quality is not the root cause of calibration failure.
The fundamental bottleneck remains the signal landscape, not the contamination injection protocol.


%% ============================================================
%% 13  DISCUSSION
%% ============================================================
\section{Discussion}

\paragraph{A diagnostic framework, not a calibration tool.}
Our results demonstrate that the conditions for successful contamination-aware calibration---reliable detection signals across multiple methods, sufficient contamination effect sizes, and stable identifiability---are not met in the SFT paraphrase regime.
Rather than viewing this as a failure, we propose reframing contamination-aware evaluation as a diagnostic problem: \emph{can we determine when calibration is feasible before attempting it?}
The identifiability boundaries (AUC~$\geq 0.80$), the Min-K++ reversal test (negative Cohen's $d$ as a red flag), and the format learning diagnostic ($\gamma$ confounded when format changes dominate content exposure) provide concrete, actionable checks.

\paragraph{Implications for the contamination detection community.}
The Min-K++ reversal finding (Section~\ref{sec:minkpp}) suggests that perplexity-based detection methods may be fundamentally unreliable for paraphrased contamination.
The reversal worsens with scale ($d$ from $-1.98$ to $-2.7$) and generalizes across model families, indicating a systematic failure mode.
Future detection methods should be validated not only for sensitivity (true positive rate) but also for \emph{directional consistency} under high-dosage contamination.

\paragraph{Why adjustment hurts.}
The counterintuitive finding that raw scores outperform all calibration methods (MAE 4.2~pp vs.\ 8.7~pp) has a simple explanation: in our SFT paraphrase setting, contamination effects are small ($\leq 6$~pp even at d100/7B) and often negative (catastrophic forgetting).
Any post-hoc adjustment adds noise proportional to detection signal error, which exceeds the contamination signal itself.
This is not a flaw of IRT specifically---linear regression shows the same pattern (MAE 8.1~pp).

\paragraph{External validity.}
Our SFT injection protocol produces paraphrased contamination, which may underestimate real-world contamination effects (reported at 8--15~pp for verbatim leakage).
The signal reliability findings and Min-K++ reversal should generalize to any setting involving non-verbatim contamination; the calibration failure results may not apply to high-effect verbatim contamination where signals are stronger.


%% ============================================================
%% 14  LIMITATIONS
%% ============================================================
\section{Limitations}

\begin{itemize}
\item \textbf{Single-family controlled experiments.} Controlled injection is limited to Qwen2.5 (with Llama cross-validation). Different tokenizers, pre-training corpora, and architectures may produce different contamination dynamics.

\item \textbf{Paraphrase vs.\ verbatim contamination.} Our protocol injects paraphrased items. Verbatim contamination may produce larger effect sizes and different signal reliability patterns.

\item \textbf{Local independence violation.} IRT assumes independent examinees; a single LLM producing all responses violates this assumption. Per-benchmark fitting and posterior predictive checks (residual correlation 0.18) partially address but do not resolve this.

\item \textbf{Unidimensional $\theta$.} We model each benchmark with a single ability dimension. MMLU's heterogeneous subject domains likely require multidimensional IRT \citep{reckase2009multidimensional} for accurate item modeling.

\item \textbf{Closed-weight infeasibility.} The framework requires embedding or logit access for effective signals, limiting applicability to open-weight models.
\end{itemize}


%% ============================================================
%% 15  CONCLUSION
%% ============================================================
\section{Conclusion}

We present a systematic empirical study of contamination signal reliability that reveals fundamental challenges for contamination-aware LLM evaluation.
Three findings have immediate implications:
Min-K++ exhibits directional reversal under paraphrased contamination (Cohen's $d$ from $-1.2$ to $-2.7$), making it actively misleading in the SFT regime;
identifiability analysis establishes that calibration requires detection signal AUC~$\geq 0.80$, a threshold unmet by three of four tested methods;
and format learning confounds contamination effect estimation on benchmarks with structured outputs.

Our ContamCal proof-of-concept honestly exposes these failure modes: it underperforms raw scores precisely because the signal landscape is too degraded.
Rather than a negative result about calibration, we view this as a positive contribution to understanding \emph{when} contamination-aware evaluation is feasible---and when the community should direct effort toward developing better detection signals before attempting calibration.

\paragraph{Future work.}
Promising directions include data attribution methods \citep{pruthi2020estimating,park2023trak} as alternative exposure signals, verbatim contamination experiments to test calibration under higher-effect regimes, and distribution-free uncertainty quantification via conformal prediction \citep{angelopoulos2023conformal}.


%% ============================================================
%% REFERENCES
%% ============================================================

\bibliographystyle{plainnat}

\begin{thebibliography}{99}

\bibitem[Angelopoulos and Bates(2023)]{angelopoulos2023conformal}
A.~N. Angelopoulos and S.~Bates.
\newblock Conformal prediction: A gentle introduction.
\newblock \emph{Foundations and Trends in Machine Learning}, 16(4):494--591, 2023.

\bibitem[Brown et~al.(2020)]{brown2020language}
T.~Brown et~al.
\newblock Language models are few-shot learners.
\newblock In \emph{NeurIPS}, 2020.

\bibitem[Cormack et~al.(2009)]{cormack2009reciprocal}
G.~V. Cormack, C.~L.~A. Clarke, and S.~B\"uttcher.
\newblock Reciprocal rank fusion outperforms condorcet and individual rank learning methods.
\newblock In \emph{SIGIR}, 2009.

\bibitem[De~Boeck and Wilson(2004)]{deboeck2004explanatory}
P.~De~Boeck and M.~Wilson.
\newblock \emph{Explanatory Item Response Models: A Generalized Linear and Nonlinear Approach}.
\newblock Springer, 2004.

\bibitem[Dekoninck et~al.(2024)]{dekoninck2024constat}
J.~Dekoninck et~al.
\newblock {ConStat}: Performance-based contamination detection in large language models.
\newblock In \emph{ICML}, 2024.

\bibitem[Deng et~al.(2024)]{deng2024investigating}
C.~Deng et~al.
\newblock Investigating data contamination in modern benchmarks for large language models.
\newblock In \emph{NAACL}, 2024.

\bibitem[Eckerly(2017)]{eckerly2017detecting}
C.~A. Eckerly.
\newblock Detecting item preknowledge and item compromise.
\newblock In \emph{Handbook of Quantitative Methods for Detecting Cheating on Tests}, pages 125--150. Routledge, 2017.

\bibitem[Guo et~al.(2017)]{guo2017calibration}
C.~Guo, G.~Pleiss, Y.~Sun, and K.~Q. Weinberger.
\newblock On calibration of modern neural networks.
\newblock In \emph{ICML}, 2017.

\bibitem[Holland and Wainer(1993)]{holland1993differential}
P.~W. Holland and H.~Wainer.
\newblock \emph{Differential Item Functioning}.
\newblock Lawrence Erlbaum Associates, 1993.

\bibitem[Jacovi et~al.(2023)]{jacovi2023stop}
A.~Jacovi et~al.
\newblock Stop uploading test data in plain text: Practical strategies for mitigating data contamination by evaluation benchmarks.
\newblock In \emph{EMNLP}, 2023.

\bibitem[Lalor et~al.(2024)]{lalor2024benchmarking}
J.~P. Lalor et~al.
\newblock Benchmarking large language models with item response theory.
\newblock In \emph{AAAI}, 2024.

\bibitem[Mislevy(1986)]{mislevy1986bayes}
R.~J. Mislevy.
\newblock Bayes modal estimation in item response models.
\newblock \emph{Psychometrika}, 51(2):177--195, 1986.

\bibitem[Oren et~al.(2024)]{oren2024proving}
Y.~Oren et~al.
\newblock Proving test set contamination in black-box language models.
\newblock In \emph{ICLR}, 2024.

\bibitem[Park et~al.(2023)]{park2023trak}
S.~M. Park et~al.
\newblock {TRAK}: Attributing model behavior at scale.
\newblock In \emph{ICML}, 2023.

\bibitem[Platt(1999)]{platt1999probabilistic}
J.~Platt.
\newblock Probabilistic outputs for support vector machines and comparisons to regularized likelihood methods.
\newblock \emph{Advances in Large Margin Classifiers}, 10(3):61--74, 1999.

\bibitem[Polo et~al.(2024)]{polo2024tinybenchmarks}
F.~M. Polo et~al.
\newblock {tinyBenchmarks}: Evaluating {LLMs} with fewer examples.
\newblock In \emph{ICML}, 2024.

\bibitem[Pruthi et~al.(2020)]{pruthi2020estimating}
G.~Pruthi et~al.
\newblock Estimating training data influence by tracing gradient descent.
\newblock In \emph{NeurIPS}, 2020.

\bibitem[Reckase(2009)]{reckase2009multidimensional}
M.~D. Reckase.
\newblock \emph{Multidimensional Item Response Theory}.
\newblock Springer, 2009.

\bibitem[Sainz et~al.(2023)]{sainz2023chatgpt}
O.~Sainz et~al.
\newblock {ChatGPT}: Jack of all trades, master of none.
\newblock \emph{arXiv preprint arXiv:2302.10724}, 2023.

\bibitem[Shi et~al.(2024)]{shi2024detecting}
W.~Shi et~al.
\newblock Detecting pretraining data from large language models.
\newblock In \emph{ICLR}, 2024.

\bibitem[Sinharay(2017)]{sinharay2017assessing}
S.~Sinharay.
\newblock Assessing item fit for unidimensional item response theory models using residuals from estimated item response functions.
\newblock \emph{Psychometrika}, 82(1):126--145, 2017.

\bibitem[Thissen et~al.(1993)]{thissen1993detection}
D.~Thissen, L.~Steinberg, and H.~Wainer.
\newblock Detection of differential item functioning using the parameters of item response models.
\newblock In \emph{Differential Item Functioning}, pages 67--113. Lawrence Erlbaum, 1993.

\bibitem[White et~al.(2024)]{white2024livebench}
C.~White et~al.
\newblock {LiveBench}: A challenging, contamination-free {LLM} benchmark.
\newblock \emph{arXiv preprint arXiv:2406.19314}, 2024.

\bibitem[Xu et~al.(2024)]{xu2024benchmarking}
R.~Xu et~al.
\newblock Benchmarking benchmark leakage in large language models.
\newblock In \emph{ICML}, 2024.

\bibitem[Yang et~al.(2023)]{yang2023rethinking}
S.~Yang et~al.
\newblock Rethinking benchmark and contamination for language models with rephrased samples.
\newblock \emph{arXiv preprint arXiv:2311.04850}, 2023.

\bibitem[Zumbo(1999)]{zumbo1999handbook}
B.~D. Zumbo.
\newblock A handbook on the theory and methods of differential item functioning ({DIF}).
\newblock Technical report, National Defense Headquarters, Ottawa, 1999.

\end{thebibliography}

\end{document}
